\documentclass[11pt]{ctexart}
\usepackage[a4paper,margin=1in]{geometry}
\usepackage{graphicx}
\usepackage{xcolor}
\usepackage{hyperref}
\definecolor{refgray}{gray}{0.45}
\title{\textbf{市场最新观点汇总}\\[0.4em]{\large 市场主线从总量复苏转向结构性分化：AI与半导体设备竞赛深化、中国供给侧再定价、美国资金虹吸效应结构裂变、欧洲财政扩张高度依赖德国。}}
\author{KC桌面 自动生成}
\date{260518}
\begin{document}
\maketitle
\section*{一页摘要}
\begin{itemize}
  \item 韩国半导体设备进口创历史新高，HBM和1c DRAM节点驱动存储器双雄的产能竞赛，测试设备和光刻设备需求结构分化。
  \item 中国4月经济数据显示工业生产与消费双双走弱，PPI超预期来自上游资源品而非终端需求，居民部门仍在去杠杆。
  \item 全球资金持续流入美国权益资产，但债券市场出现外流，新兴市场资金流入断崖式回落。
  \item 欧元区2026年财政扩张几乎完全由德国驱动，法国、意大利、西班牙规划紧缩，财政脉冲高度依赖德国执行。
  \item AI对竞争格局的影响存在两面性：可能加剧头部企业集中度，也可能打破现有垄断，历史数据显示技术冲击后集中度往往加速。
  \item 中国衍生品新规通过提高资本门槛加速行业集中度提升，头部券商竞争优势扩大。
  \item 汽车股alpha来自非汽车叙事，储能和AI数据中心成为估值重构的关键主题。
\end{itemize}
\section{AI/半导体与设备投资}
\textbf{韩国半导体设备进口数据确认存储器双雄在HBM和1c节点上的不可逆产能竞赛，设备需求结构正在被重塑。}

\begin{itemize}
  \item 2026年4月韩国半导体设备进口总额同比增长57\%至34.5亿美元，创历史新高，尽管环比下降9\%，但3个月移动平均仍在上升。
  \item 设备进口与资本开支的季度背离（1Q26合并资本开支环比下降）并非需求减弱，而是2025年四季度基础设施投资前置的结果，两家公司均指引2026年资本开支大幅增长。
  \item 来自荷兰的光刻设备进口环比回落28\%，但仍是历史上第二高的季度首月数据，同比暴增约180\%，DRAM扩产和1c节点快速导入是主要驱动力。
  \item 韩国测试设备进口（来自日本+马来西亚）4月环比-21\%，但同比+49\%，回归模型暗示爱德万6月季度韩国收入可能环比+70\%，远超市场预期的+5\%。
  \item 日本设备进口（CVD、刻蚀、清洗等）4月环比-19\%，东京电子短期承压，但长期与韩国资本开支周期高度相关。
\end{itemize}
\begin{figure}[htbp]
\centering
\includegraphics[width=0.88\linewidth]{figures/fig_001.jpg}
\caption{Exhibit 1 - EXHIBIT 1: April South Korea imports for SPE was \textbackslash{}\$3.4bn, -9\% MoM.}
\end{figure}
\begin{figure}[htbp]
\centering
\includegraphics[width=0.88\linewidth]{figures/fig_003.jpg}
\caption{EXHIBIT 3 - EXHIBIT 3: South Korea SPE monthly import data has good directional correlation with Samsung and SK hynix's quarterly capex.}
\end{figure}
\begin{figure}[htbp]
\centering
\includegraphics[width=0.88\linewidth]{figures/fig_004.jpg}
\caption{EXHIBIT 4 - EXHIBIT 4: 1QCY26 capex came down QoQ for both Samsung and SK hynix due to seasonality and also as some front loaded infrastructure investments in 4Q25. Capex (KRW B)}
\end{figure}
\begin{figure}[htbp]
\centering
\includegraphics[width=0.88\linewidth]{figures/fig_005.jpg}
\caption{Exhibit 6 - EXHIBIT 5: April tester imports from Japan and Malaysia collectively was -21\% MoM.}
\end{figure}
\begin{figure}[htbp]
\centering
\includegraphics[width=0.88\linewidth]{figures/fig_006.jpg}
\caption{EXHIBIT 6 - EXHIBIT 6: South Korea tester imports data shows good directional correlation with Advantest's South Korea sales.}
\end{figure}
\begin{figure}[htbp]
\centering
\includegraphics[width=0.88\linewidth]{figures/fig_007.jpg}
\caption{EXHIBIT 8 - EXHIBIT 8: Advantest's Quarterly Korean sales shows good correlation with Korean monthly imports. Advantest Quarterly KR Sales vs. KR 1M Import: 1QCY16-1QCY26}
\end{figure}
\begin{figure}[htbp]
\centering
\includegraphics[width=0.88\linewidth]{figures/fig_011.jpg}
\caption{Exhibit 13 - EXHIBIT 13: South Korea monthly SPE imports from Netherlands shows good directional correlation with ASML's South Korea revenue divided by 3. April Litho Imports of EUR 723mn declined 28\% MoM but up \textbackslash{}\textasciitilde{}180\% YoY.}
\end{figure}
\begin{figure}[htbp]
\centering
\includegraphics[width=0.88\linewidth]{figures/fig_012.jpg}
\caption{EXHIBIT 14 - EXHIBIT 14: Regressing ASML Korea's quarterly system sales against 1 month + 1M lag of import data, we obtain an R \$\textasciicircum{}\{2\}\$ of 75\% and estimate that Q2 ASML Korea revenue will reach EUR 2.1bn, down 26\% QoQ.}
\end{figure}
{\small\color{refgray}\textbf{References:} [R001] 韩国设备进口数据揭示的真正信号：HBM和1c节点正在重塑半导体设备的需求结构}

\section{中国宏观与利率}
\textbf{中国经济核心矛盾正从需求端能否复苏转向供给侧如何被重新定价，PPI超预期、居民去杠杆与中美贸易结构性矛盾并存。}

\begin{itemize}
  \item 4月PPI同比跃升至2.8\%，远超市场预期的1.8\%，跳升来自上游资源品价格的结构性推升（石油开采+28.6\%，有色金属开采+38.9\%），而非终端需求全面复苏。
  \item 4月工业增加值同比增速从3月的5.7\%骤降至4.1\%，出口增长中约一半来自芯片和电子产品价格飙升，实际出口量增长有限。
  \item 4月社零同比仅+0.2\%，扣除通胀后实际零售增速转负（-1.0\%），汽车销售同比-15.3\%，家电-15.0\%，以旧换新政策退坡效应显现。
  \item 居民部门仍在加速去杠杆，4月居民人民币贷款净增量持续为负，家庭债务/GDP从2024年一季度的62.3\%降至去年底的59.4\%。
  \item 中美首脑会晤后结构性矛盾未消解，短期风险后移，但双方侧重完全不同——中方强调战略稳定，美方聚焦购买美国商品和霍尔木兹海峡通航。
\end{itemize}
{\small\color{refgray}\textbf{References:} [R003] 市场真正低估的不是需求，而是供给侧的再定价; [R006] 四月经济数据证实：市场低估的不是反弹结束，而是结构性分化的深度}

\section{地产与消费}
\textbf{新房价格跌幅收窄是真实的边际改善，但不足以支撑反转判断，二手房市场以5\%-15\%的年度跌幅持续定价，新房价格信号已失去代表性。}

\begin{itemize}
  \item 4月70城新房加权均价环比年化跌幅从3月的4.0\%收窄至3.0\%，一线城市继续环比上涨，上海领涨全国（环比年化+2.0\%）。
  \item 同比数据并不乐观，全国加权均价同比跌幅从3月的-3.5\%扩大至-3.6\%，真正实现同比上涨的城市仅上海(+3.7\%)和杭州(+2.3\%)。
  \item 二线城市环比跌幅从-4.6\%收窄至-2.8\%，三线城市从-5.0\%收窄至-4.6\%，收窄是事实但绝对值仍然很高，市场分化在加剧。
  \item 高频数据显示30城新房交易量在4月和5月初基本持平去年同期，主要城市库存去化周期从29.3个月降至28.9个月，但新房价格环比上涨的城市数量在减少。
  \item 二手房市场调整更深，统计局和第三方平台数据显示过去一年跌幅在5\%-15\%之间，新房与二手房价格信号正在走向两个世界。
\end{itemize}
\begin{figure}[htbp]
\centering
\includegraphics[width=0.88\linewidth]{figures/fig_014.jpg}
\caption{Exhibit 1 - Exhibit 1: In the primary market, the sequential decline in 70-city weighted average property price eased to \$3.0\textbackslash{}\%\$ mom annualized in April}
\end{figure}
\begin{figure}[htbp]
\centering
\includegraphics[width=0.88\linewidth]{figures/fig_015.jpg}
\caption{Exhibit 2 - Exhibit 2: The share of cities that experienced sequentially higher property prices edged down in secondary markets in April}
\end{figure}
\begin{figure}[htbp]
\centering
\includegraphics[width=0.88\linewidth]{figures/fig_016.jpg}
\caption{Exhibit 3 - Exhibit 3: NBS 70-city secondary home prices continued to decline in April}
\end{figure}
{\small\color{refgray}\textbf{References:} [R002] 新房价格跌幅收窄，但真正的市场底不在统计局数据里}

\section{全球资金流向与区域市场}
\textbf{美国资产的虹吸效应并未减弱，但结构正在裂变——权益与债券偏好反转，新兴市场资金流入出现断崖式回落。}

\begin{itemize}
  \item 5月8日至14日当周，外国资金净流入美国相关基金21亿美元，权益类基金贡献24亿美元净流入，债券类基金净流出2.92亿美元。
  \item 5月至今美国权益类基金累计流入46亿美元，远超4月全月的31亿美元，美国权益资产吸引力在加速强化。
  \item 新兴市场ETF同期净流出9700万美元，逆转前一周12亿美元净流入，债券基金流出2.17亿美元是主要拖累。
  \item 韩国散户继续净卖出美国资产7200万美元，但5月至今累计净卖出1400万美元，较4月的9.1亿净卖出明显收窄。
  \item 台湾投资者买入美元债券ETF 1.03亿美元，5月至今累计净买入3.23亿，已扭转4月净卖出3.42亿的局面。
\end{itemize}
\begin{figure}[htbp]
\centering
\includegraphics[width=0.88\linewidth]{figures/fig_028.jpg}
\caption{Figure 3 - - Foreign inflows into US-focused equity ETFs and mutual funds \$\textasciicircum{}\{[1]\}\$ totaled USD2.4bn between 8 and 14 May (Figure 3), continuing from inflows of USD2.2bn in the previous week. MTD May, inflows into these funds totaled USD4.6bn surpassing total inflows of U}
\end{figure}
{\small\color{refgray}\textbf{References:} [R005] 全球资金流向正在发出一个被忽视的信号：美国资产的“虹吸效应”并未减弱，但结构正在裂变}

\section{欧洲财政与区域市场}
\textbf{欧元区2026年财政扩张几乎完全由德国一国驱动，法国、意大利、西班牙规划紧缩，财政脉冲高度依赖德国执行。}

\begin{itemize}
  \item 2026年欧元区整体赤字率将从2.9\%上升至3.5\%，但经周期调整的初级财政赤字仅扩大约0.2个百分点，远低于市场部分预期。
  \item 德国一国贡献了0.5个百分点的宽松力度，而法国、意大利、西班牙各自贡献0.1个百分点的紧缩，恰好抵消德国以外的全部扩张。
  \item 21个欧元区国家中，11个赤字在扩大，9个在缩小，1个持平，德国赤字扩大1.6个百分点至4.3\%，爱沙尼亚扩大2.3个百分点至4.3\%。
  \item 如果剔除德国，欧元区2026年财政立场实际上是略微紧缩的，市场对欧洲财政扩张的定价隐含了德国必然执行的前提假设。
  \item 德国联邦支出在第一季度已由国防开支驱动增长，但预算计划只是计划，执行才是现实，德国能否兑现庞大财政扩张承诺是关键变量。
\end{itemize}
\begin{figure}[htbp]
\centering
\includegraphics[width=0.88\linewidth]{figures/fig_031.jpg}
\caption{Figure 1 - Figure 1: Eurozone headline budget balance, \% of GDP}
\end{figure}
\begin{figure}[htbp]
\centering
\includegraphics[width=0.88\linewidth]{figures/fig_032.jpg}
\caption{Figure 2 - Figure 2: Eurozone revenues and expenditures, \% of GDP}
\end{figure}
\begin{figure}[htbp]
\centering
\includegraphics[width=0.88\linewidth]{figures/fig_033.jpg}
\caption{Figure 3 - Figure 3: Headline budget deficits, 2026, \% of GDP}
\end{figure}
\begin{figure}[htbp]
\centering
\includegraphics[width=0.88\linewidth]{figures/fig_034.jpg}
\caption{Figure 4 - Figure 4: Eurozone, level of cyclically adjusted primary balance, \% of GDP}
\end{figure}
\begin{figure}[htbp]
\centering
\includegraphics[width=0.88\linewidth]{figures/fig_035.jpg}
\caption{Figure 5 - Figure 5: Eurozone, change in cyclically adjusted primary balance, \% of GDP}
\end{figure}
\begin{figure}[htbp]
\centering
\includegraphics[width=0.88\linewidth]{figures/fig_037.jpg}
\caption{Figure 7 - Figure 7: Contributions to 2026 Eurozone fiscal impulse, \% of GDP}
\end{figure}
\begin{figure}[htbp]
\centering
\includegraphics[width=0.88\linewidth]{figures/fig_040.jpg}
\caption{Figure 11 - Figure 11: Germany budget deficit and CAPB, \% of GDP}
\end{figure}
\begin{figure}[htbp]
\centering
\includegraphics[width=0.88\linewidth]{figures/fig_041.jpg}
\caption{Figure 12 - Figure 12: Italy budget deficit and CAPB, \% of GDP}
\end{figure}
\begin{figure}[htbp]
\centering
\includegraphics[width=0.88\linewidth]{figures/fig_042.jpg}
\caption{Figure 13 - Figure 13: France budget deficit and CAPB, \% of GDP}
\end{figure}
\begin{figure}[htbp]
\centering
\includegraphics[width=0.88\linewidth]{figures/fig_043.jpg}
\caption{Figure 14 - Figure 14: Spain budget deficit and CAPB, \% of GDP}
\end{figure}
{\small\color{refgray}\textbf{References:} [R008] 欧元区2026年财政的真实底色：扩张全靠德国，紧缩才是多数国家的选择}

\section{AI与竞争格局}
\textbf{AI究竟是加剧垄断还是打破垄断，将决定未来十年资产定价的核心逻辑，而当前企业集中度和利润率均处于历史高位。}

\begin{itemize}
  \item 美国企业集中度自1930年代以来持续上升，头部1\%企业的销售额占比从1960年代的约60\%攀升至近年来的约80\%，这一趋势在多个发达经济体同样存在。
  \item 全球化与反垄断执法弱化并非集中度上升的主因，跨国面板数据检验显示贸易开放度与集中度上升之间缺乏显著关联，不同反垄断执法周期内集中度上升速度大致相同。
  \item 技术驱动是集中度上升最站得住脚的解释，新技术让少数公司靠规模效应吃掉更多市场，技术冲击后集中度往往加速。
  \item 过去40年美国企业利润率大幅上升，报告估算2000年以来集中度上升解释了利润率增长约三分之一。
  \item AI可能打破现有格局让竞争更激烈，也可能让AI用得好的公司跑得更远，进一步提高集中度，AI暴露度高的行业集中度和利润率均高于平均水平。
\end{itemize}
\begin{figure}[htbp]
\centering
\includegraphics[width=0.88\linewidth]{figures/fig_017.jpg}
\caption{Exhibit 1 - Exhibit 1: Corporate Profits and Concentration Are at Record Highs in the US}
\end{figure}
\begin{figure}[htbp]
\centering
\includegraphics[width=0.88\linewidth]{figures/fig_018.jpg}
\caption{Exhibit 2 - Exhibit 2: Corporate Concentration Has Risen Steadily, Both Across Advanced Economies... Sales Share of the Top 1\% Firms by Sales Across Advanced Economies*}
\end{figure}
\begin{figure}[htbp]
\centering
\includegraphics[width=0.88\linewidth]{figures/fig_019.jpg}
\caption{Exhibit 3 - Exhibit 3: ...And US Industries Asset Share of the Top 1\% Firms by Assets*}
\end{figure}
\begin{figure}[htbp]
\centering
\includegraphics[width=0.88\linewidth]{figures/fig_020.jpg}
\caption{Exhibit 4 - Exhibit 4: Corporate Concentration Has Risen Across Different Antitrust Regimes}
\end{figure}
\begin{figure}[htbp]
\centering
\includegraphics[width=0.88\linewidth]{figures/fig_021.jpg}
\caption{Exhibit 5 - Exhibit 5: Corporate Concentration Has Risen Faster During Periods of Rapid Technological Change}
\end{figure}
\begin{figure}[htbp]
\centering
\includegraphics[width=0.88\linewidth]{figures/fig_022.jpg}
\caption{Exhibit 6 - Exhibit 6: Technological Change Has Widened Productivity Dispersion Across Firms, Particularly in Technology-Related Industries, Increasing the Gap Between Frontier Firms and the Rest}
\end{figure}
\begin{figure}[htbp]
\centering
\includegraphics[width=0.88\linewidth]{figures/fig_023.jpg}
\caption{Exhibit 7 - Exhibit 7: Markups Charged by US Firms Have Risen Since the 1980s...}
\end{figure}
\begin{figure}[htbp]
\centering
\includegraphics[width=0.88\linewidth]{figures/fig_024.jpg}
\caption{Exhibit 8 - Exhibit 8: ...Partly Reflecting Lower Consumer Price Sensitivity as Rising Incomes Lead Households to Search Less Across Sellers Change in US Consumer Sensitivity to Prices Relative to 2006 Estimated at the Product-Level*}
\end{figure}
\begin{figure}[htbp]
\centering
\includegraphics[width=0.88\linewidth]{figures/fig_025.jpg}
\caption{Exhibit 9 - Exhibit 9: Rising Concentration Accounts for Roughly One-Third of the Increase in Profit Margins Since 2000 Change in Corporate Profit Margins, 2000-24*}
\end{figure}
\begin{figure}[htbp]
\centering
\includegraphics[width=0.88\linewidth]{figures/fig_027.jpg}
\caption{Exhibit 11 - Exhibit 11: AI-Exposed Industries Are Somewhat More Concentrated and Operate with Higher Profit Margins Relationship Between Industry-Level Concentration and AI Exposure}
\end{figure}
{\small\color{refgray}\textbf{References:} [R004] 市场真正低估的不是AI需求，而是AI对竞争格局的再定价}

\section{中国衍生品监管与券商格局}
\textbf{衍生品新规不是市场的收缩，而是行业集中度加速提升的起点，头部券商竞争优势扩大。}

\begin{itemize}
  \item 新规要求开展衍生品交易的证券期货公司近六个月净资本不低于5000万元人民币，资本实力成为入场券而非加分项。
  \item 新规扩大交易限制，禁止上市公司、大股东、高管等利用衍生品交易自家股票或类似股权证券，明确衍生品应服务套期保值，抑制过度投机。
  \item 不合规产品要在6个月过渡期内逐步清退，资本实力弱的小机构可能退出市场，但长期为更规范、更有序的杠杆使用打开空间。
  \item CITIC、CICC、HTSC、GTJA四家头部券商2025年权益衍生品名义本金占所有上市券商的近60\%（2024年约52\%），集中度已在提升。
  \item 新规通过提高资本要求、扩大禁止交易范围、强化分类与隔离要求，系统性地抬高衍生品业务运营成本，小机构被迫退出，龙头券商凭借合规成本和产品供给能力双重优势扩大份额。
\end{itemize}
\begin{figure}[htbp]
\centering
\includegraphics[width=0.88\linewidth]{figures/fig_029.jpg}
\caption{Figure 1 - Figure 1: Leading brokers' share in notional principal of equity derivatives expanded Notional principal of equity derivatives}
\end{figure}
\begin{figure}[htbp]
\centering
\includegraphics[width=0.88\linewidth]{figures/fig_030.jpg}
\caption{Figure 2 - Figure 2: Commodities derivatives dominating exchange-traded in China Exchange-traded derivatives}
\end{figure}
{\small\color{refgray}\textbf{References:} [R007] 中国衍生品新规的真正信号：不是收紧，是龙头加速分化的发令枪}

\section{汽车产业链与跨界叙事}
\textbf{在传统汽车产量持续承压的背景下，市场真正奖励的是能将技术能力迁移到AI数据中心和储能等非汽车领域的公司。}

\begin{itemize}
  \item S\&P Global Mobility将2026年全球汽车产量同比降幅从-1.8\%进一步扩大至-2.4\%，但下修集中在2026年二至四季度，一季度实际上调了+1.6\%。
  \item 产量下修中约41\%来自中国，北美仅占5\%（约4.2万辆），对北美供应商比表面数字更友好。
  \item 福特凭借储能业务（BESS）的想象空间单周上涨9\%，BorgWarner和Aptiv等供应商也因非汽车主题获得资金追捧。
  \item 2027年预期也被下调-1.3\%，但只是同比+0.8\%，研报推测S\&P一次性调得较多是为了避免慢慢砍的痛苦，若宏观担忧缓解未来有上调空间。
  \item 汽车AI交易并未消亡，但涨速过快，福特与CATL的技术合作、PTC等成为市场关注焦点。
\end{itemize}
\begin{figure}[htbp]
\centering
\includegraphics[width=0.88\linewidth]{figures/fig_046.jpg}
\caption{Figure 2 - Figure 2: US Automotive coverage stock performance, past week}
\end{figure}
\begin{figure}[htbp]
\centering
\includegraphics[width=0.88\linewidth]{figures/fig_047.jpg}
\caption{Figure 3 - Figure 3: US Automotive coverage stock performance, last 3 months}
\end{figure}
\begin{figure}[htbp]
\centering
\includegraphics[width=0.88\linewidth]{figures/fig_048.jpg}
\caption{Figure 4 - Figure 4: US Automotive coverage stock performance, last 12 months Stock performance, 12 months}
\end{figure}
\begin{figure}[htbp]
\centering
\includegraphics[width=0.88\linewidth]{figures/fig_051.jpg}
\caption{Figure 12 - Figure 12: US Monthly SAAR Trend}
\end{figure}
\begin{figure}[htbp]
\centering
\includegraphics[width=0.88\linewidth]{figures/fig_061.jpg}
\caption{Figure 22 - Figure 22: North America LVP - S\&P Forecast North America LVP}
\end{figure}
\begin{figure}[htbp]
\centering
\includegraphics[width=0.88\linewidth]{figures/fig_062.jpg}
\caption{Figure 23 - Figure 23: Europe LVP - S\&P Forecast Europe LVP}
\end{figure}
\begin{figure}[htbp]
\centering
\includegraphics[width=0.88\linewidth]{figures/fig_063.jpg}
\caption{Figure 24 - Figure 24: China LVP - S\&P Forecast China LVP}
\end{figure}
\begin{figure}[htbp]
\centering
\includegraphics[width=0.88\linewidth]{figures/fig_064.jpg}
\caption{Figure 25 - Figure 25: Global LVP - S\&P Forecast Global LVP}
\end{figure}
{\small\color{refgray}\textbf{References:} [R009] 汽车股真正的alpha不在车本身，而在“非汽车”叙事能否重构估值}

\section*{结语}
综合来看，当前市场主线已从总量复苏转向结构性分化：AI与半导体设备竞赛深化、中国供给侧再定价、美国资金虹吸效应结构裂变、欧洲财政扩张高度依赖德国。理解这些分化的具体机制，比争论总量数据何时见底更具实际意义。
\vfill
{\small\color{refgray}Personal reading notes and learning share only. Not investment advice.}
\end{document}
